\documentclass[11pt]{article}
\usepackage[T1]{fontenc}
\usepackage[utf8]{inputenc}
\usepackage{lmodern}
\usepackage{amsmath,amssymb,amsthm}
\usepackage{microtype}
\usepackage{geometry}
\usepackage{enumitem}

\newtheorem{theorem}{Theorem}
\newtheorem{proposition}{Proposition}
\newtheorem{lemma}{Lemma}
\newtheorem{corollary}{Corollary}
\newtheorem{conjecture}{Conjecture}
\newtheorem{remark}{Remark}

\newtheorem{defn}[theorem]{Definition}
\newcommand{\VCdim}{\operatorname{VCdim}}

\newcommand{\cm}{\operatorname{cm}}


\title{A note on Open AI's "A counterexample to Hadwiger’s conjecture"}
\author{Pei Wu}
\date{\today}

\begin{document}
	
	\maketitle
	
	\section{Summary}
	The manuscript constructs a triangle-free graph $G$ whose complement $H=\bar G$ satisfies $\alpha(H)\le2$ and $\cm(H)<m/100$, where $m$ is the number of vertices. The connected matching number is defined as follows:
	
	\begin{defn}
		A matching $M$ in a graph $H$ is called a \emph{connected matching} if every two distinct edges of $M$ are joined by an edge of $H$: for all distinct $e,f\in M$, there exist $u\in V(e)$ and $v\in V(f)$ such that $uv\in E(H)$. The \emph{connected matching number} is
	\[
	\cm(H)=\max\bigl\{|M|: M\text{ is a connected matching in }H\bigr\}.
	\]
	\end{defn}
	
	The complement graph provides a counterexample to Hadwiger's conjecture, which asserts that $\chi(H)\le h(H)$, where $h(H)$ is the largest integer $t$ such that $H$ contains a $K_t$ minor. Indeed, Proposition~3.5 of the paper (proved in Appendix~\ref{app:minor-bound}) gives
	\[
	h(H)\le \frac{|V(H)|+4\cm(H)+2}{3}.
	\]
	Applying this to $H=\bar G$, and using $\alpha(\bar G)\le 2$ so that every color class has at most two vertices, we obtain, for $m$ large enough,
	\[
	h(\bar G)
	\le \frac{m+4\cm(\bar G)+2}{3}
	< \frac{26m}{75}+\frac{2}{3}
	< \frac{m}{2}
	\le \chi(\bar G).
	\]
	Thus $h(\bar G)<\chi(\bar G)$, contrary to Hadwiger's conjecture.
\section{The construction}
\subsection{Triangle-free graphs from linear algebra}\label{sec:2.1}
We first describe the general linear-algebraic scheme and then specify the data used in the manuscript. All vector spaces, ranks, and matrix operations are over $\mathbb F_2$.

Let $V,W$ be finite-dimensional vector spaces. Fix $a\in V^*$ and a symmetric bilinear form $T:V\times V\to\mathbb F_2$, and put
\[
r(\lambda)=a+T(\lambda,\cdot)\in V^*.
\]
Let $\Omega$ be the set of vertices. For each $o\in\Omega$, assign an injective linear map $U_o:V\to W$ and a linear functional $u_o\in W^*$. Then we can define the associated graph:

\begin{defn}
Two vertices $i,j\in\Omega$ are adjacent if there exist $\lambda_i,\lambda_j\in V$ such that
\begin{align}
U_i\lambda_i&=U_j\lambda_j,\label{eq:shared}\\
u_j\circ U_i&=r(\lambda_i),\qquad u_i\circ U_j=r(\lambda_j),\label{eq:cross}\\
a(\lambda_i)+a(\lambda_j)&=1.\label{eq:parity}
\end{align}
Equation~\eqref{eq:cross} is an equality in $V^*$.
\end{defn}

\begin{lemma}
The associated graph is loopless and triangle-free.
\end{lemma}
\begin{proof}
	Loopless is guaranteed by injectivity of $U$ and (\ref{eq:parity}).
Suppose $i,j,k$ is a triangle. Write $\lambda_{ij}$ for the $\lambda_i$ at $i$ for the pair $ij$.  Evaluate (\ref{eq:cross}) on $\lambda_{ik}$:
\[
u_j(U_i \lambda_{ik})=a(\lambda_{ik})+T(\lambda_{ij},\lambda_{ik}).
\]
Then sum six equations by permuting $i,j,k$. For fixed $j$, the two terms on the left cancel because $U_i \lambda_{ik}=U_k \lambda_{ki}$. For fixed $i$, the two bilinear terms cancel by symmetry of $T$. The remaining right-hand side is
\[
\sum_{\{i,j\}}\bigl(a(\lambda_{ij})+a(\lambda_{ji})\bigr)=1+1+1=1,
\]
whereas the left-hand side is zero, a contradiction.
\end{proof}
\subsection{Construction of $U$ and $u$}\label{sec:frames}
Write $\operatorname{Mat}_{m,n}(\mathbb F_2)$ for the space of $m\times n$ binary matrices, and $\operatorname{Mat}_n(\mathbb F_2)=\operatorname{Mat}_{n,n}(\mathbb F_2)$. Fix a finite index set $\mathcal E$, positive integers $N,d_B,h$ with $N\ge2(d_B+h)$, and a matrix $E\in\operatorname{Mat}_{d_B}(\mathbb F_2)$, to be specified later. Take
\[
V\subseteq\bigoplus_{e\in\mathcal E}\operatorname{Mat}_{d_B}(\mathbb F_2)
\quad\text{a linear subspace},\qquad
W=\bigoplus_{e\in\mathcal E}\operatorname{Mat}_N(\mathbb F_2).
\]
For $x\in V$, write $x=(x_e)_{e\in\mathcal E}$, so $x_e\in\operatorname{Mat}_{d_B}(\mathbb F_2)$.

At each component $e$, consider quadruples
\[
P_e,Q_e\in\operatorname{Mat}_{N,d_B}(\mathbb F_2),\qquad
X_e,Y_e\in\operatorname{Mat}_{N,h}(\mathbb F_2)
\]
satisfying
\begin{align*}
	[P_e\;X_e]&\text{ and }[Q_e\;Y_e]\text{ have full column rank},\\
	P_e^TQ_e&=E,\\
	X_e^TQ_e&=0,\qquad P_e^TY_e=0.
\end{align*}
The condition $N\ge2(d_B+h)$ ensures there is enough admissible quadruples. Let $\mathcal{O}$ be the set of tuples $o=(P_{o,e},Q_{o,e},X_{o,e},Y_{o,e})_{e\in\mathcal E}$ of admissible quadruples. Define a probability law on $\mathcal{O}$ by sampling uniformly at each component, independently across components. 

For matrices of the same size, define $\langle A,B\rangle=\operatorname{Tr}(A^TB)$. Each $o\in\mathcal O$ determines $U_o:V\to W$ and $u_o\in W^*$ by
\[
U_o(x)=(P_{o,e}x_eQ_{o,e}^T)_{e\in\mathcal E},\qquad
u_o(\xi)=\sum_{e\in\mathcal E}\langle Y_{o,e}X_{o,e}^T,\xi_e\rangle.
\]
for $x\in V$ and $\xi\in W$. Each image component $P_{o,e}x_eQ_{o,e}^T$ belongs to $\operatorname{Mat}_N(\mathbb F_2)$. Left inverses of $P_{o,e}$ and $Q_{o,e}$ recover $x_e$ from that component, so $U_o$ is injective. Thus these maps provide the vertex data required in Section~\ref{sec:2.1}; the particular subspace $V$ and the forms $a,T$ will be specified below.

Moreover,
\[
u_o(U_o(x))=\sum_{e\in\mathcal E}
\operatorname{Tr}\bigl(X_{o,e}Y_{o,e}^TP_{o,e}x_eQ_{o,e}^T\bigr)=0,
\]
since $Y_{o,e}^TP_{o,e}=0$.
\subsubsection{A model for uniformly generating quadruples} 
Set $k=d_B+h$, and combine the columns into
\[
A=[P\;X],\qquad B=[Q\;Y]\in\operatorname{Mat}_{N,k}(\mathbb F_2).
\]
The constraints say that $A,B$ have full column rank and
\[
A^TB=H_C,\qquad
H_C=\begin{pmatrix}E&0\\0&C\end{pmatrix},\qquad
C\in\operatorname{Mat}_h(\mathbb F_2).
\]
Here $C=X^TY$ is unrestricted. The letter $B$ in this subsection denotes the combined minus frame, not the coefficient space introduced later.

First fix a full-column-rank $A$. Choose a right inverse $D$ of $A^T$ and a matrix $K$ whose columns form a basis of $\ker A^T$:
\[
D\in\operatorname{Mat}_{N,k}(\mathbb F_2),\qquad
K\in\operatorname{Mat}_{N,N-k}(\mathbb F_2),\qquad
A^TD=I_k,\quad A^TK=0.
\]
Both can be obtained by Gaussian elimination, using a fixed rule for each $A$. Every matrix satisfying the linear constraints has a unique expression
\[
\boxed{B=DH_C+KZ},\qquad
C\in\operatorname{Mat}_h(\mathbb F_2),\quad
Z\in\operatorname{Mat}_{N-k,k}(\mathbb F_2).
\]
Indeed, $A^TB=H_C$ and the remaining columns lie in $\ker A^T$. Since $[D\;K]$ is invertible,
\[
\operatorname{rank}B=k
\quad\Longleftrightarrow\quad
\operatorname{rank}\begin{pmatrix}H_C\\Z\end{pmatrix}=k
\quad\Longleftrightarrow\quad
Z|_{\ker H_C}\text{ is injective}.
\]
Thus the linear equations have disappeared, leaving only an injectivity condition on the directions not detected by $H_C$.

To sample uniformly, use the following procedure:
\begin{enumerate}
\item Sample $A$ uniformly among full-column-rank matrices in $\operatorname{Mat}_{N,k}(\mathbb F_2)$, and compute $D,K$.
\item Sample $C$ and $Z$ with independent uniform binary entries. If the stacked matrix above does not have rank $k$, resample both $C$ and $Z$, keeping $A$ fixed.
\item Set $B=DH_C+KZ$ and split $A=[P\;X]$ and $B=[Q\;Y]$ at column $d_B$.
\end{enumerate}
For each $A$, the correspondence between accepted pairs $(C,Z)$ and admissible minus frames $B$ is bijective. Moreover, the acceptance condition depends only on $C,Z,E$, not on $A$, so every full-rank $A$ has the same number of admissible completions. Consequently the resulting quadruple is uniform among all admissible quadruples. Repeat this procedure independently over components.

The rejection step has a uniformly bounded expected cost. For fixed $C$, let $t=\dim\ker H_C$. The restriction of a uniform $Z$ to this kernel is a uniform $(N-k)\times t$ matrix, so its probability of being injective is
\[
\prod_{j=0}^{t-1}\bigl(1-2^{j-(N-k)}\bigr),
\]
with the empty product equal to one. Since $t\le k\le N-k$, this probability is at least $1/4$: it is bounded below by $\prod_{r=1}^{\infty}(1-2^{-r})$, whose first two factors give $3/8$ and whose remaining product is at least $1-\sum_{r=3}^{\infty}2^{-r}=3/4$. Thus fewer than four trials are needed on average. Importantly, fixing a uniformly sampled $C$ and resampling only $Z$ would generally give a different law, because the acceptance probability depends on $\operatorname{rank}C$.

\subsection{Parameters}
The only asymptotic parameter is $n$. The manuscript fixes
\[
C_0=1000,\qquad g=10^9+1,\qquad D=4C_0g,\qquad M=2^{1000}.
\]
The constants $b,j_*,J,r_0,M_0$ are chosen in the dependency order of Appendix~A of the manuscript, before $n$ tends to infinity. Their numerical values are not specified. Set
\[
h=1000r_0,\qquad N=M_0n,\qquad m=2^{C_0gN}.
\]
Here $M$ is a density bound used in the distribution theorem, whereas $M_0$ scales the ambient dimension. Choose $M_0$ sufficiently large to satisfy the dimension and counting requirements, including $N\ge2(d_B+h)$ for all sufficiently large $n$. We take $n\ge2r_0$ and $j_*\ge1$.

\subsection{Boolean moment spaces}

Let
\[
I=\mathbb Z/g\mathbb Z,\quad \mathcal E=\binom I2,\quad
I_d=\{d+1,\ldots,d+(g-1)/2\}\subset I.
\]
Elements of $I$ are \emph{tags} and elements of $\mathcal E$ are \emph{components}. A \emph{Boolean point} consists of a \emph{selector} $s\in\mathbb F_2^b$ and \emph{base coordinates} $z$ in $g^2+3$ blocks of $n$ bits:
\[
z=\bigl((z_{O_{d,t}})_{d,t\in I},\ z_S,\ z_{\#},\ z_Z\bigr)
\in(\mathbb F_2^n)^{g^2+3}.
\]
For each block label $A\in\{O_{d,t}:d,t\in I\}\cup\{S,\#,Z\}$, write
\[
z_A=(z_{A,1},\ldots,z_{A,n})\in\mathbb F_2^n.
\]
Thus $z$ has $g^2+3$ blocks, each containing $n$ bits; $z_{A,k}$ denotes the $k$th bit in block $A$.
Write
\[
p_* =\sum_{j=0}^{j_*}\binom bj,\qquad
B=\mathbb F_2^{p_*}\otimes\mathbb F_2^{1+(g^2+3)n},\qquad d_B=\dim B.
\]
Let $p(s)\in\mathbb F_2^{p_*}$ be the bit vector listing the evaluations at $s$ of all squarefree selector monomials of degree at most $j_*$, in a fixed order. Define the evaluation vector and its point moment by
\[
v(s,z)=p(s)\otimes(1,z),\qquad w(s,z)=v(s,z)v(s,z)^T\in\operatorname{Mat}_{d_B}(\mathbb F_2).
\]
Identify $B\otimes B$ with $\operatorname{Mat}_{d_B}(\mathbb F_2)$ using the fixed coefficient basis. Let $\mathcal M\subseteq\operatorname{Mat}_{d_B}(\mathbb F_2)$ be the span of all point moments. For $\ell\in I$, let $\mathcal M_\ell$ be the span of those point moments whose base values satisfy $O_{d,t}=0$ whenever $\ell\notin I_d$. Let $\mathcal M_*$ be the span obtained by setting all $O$-blocks to zero.

Define the linear cut-profile map
\[
\begin{aligned}
f:\bigoplus_{\ell\in I}\mathcal M_\ell&\longrightarrow\bigoplus_{e\in\mathcal E}\mathcal M,\\
(w_\ell)_{\ell\in I}&\longmapsto
\bigl(w_\ell+w_{\ell'}\bigr)_{\{\ell,\ell'\}\in\mathcal E}.
\end{aligned}
\]
Take the witness space in Sections~\ref{sec:2.1} and~\ref{sec:frames} to be $V=\operatorname{im}f$. A tuple $(w_\ell)_{\ell\in I}$ is a representation of $x\in V$ if $f((w_\ell)_{\ell\in I})=x$. By Lemma~2.3 of the manuscript,
\[
\ker f=\{(w_*)_{\ell\in I}:w_*\in \mathcal M_*\},
\]
that is, two representations of the same $x$ differ by a constant tuple with value in $\mathcal M_*$.

\subsection{The functional $a$, the bilinear form $T$, and the matrix $E$}
The witness space is now fixed as $V=\operatorname{im}f$, and the ambient space remains $W=\bigoplus_{e\in\mathcal E}\operatorname{Mat}_N(\mathbb F_2)$ from Section~\ref{sec:frames}. We next specify $a\in V^*$, a symmetric bilinear form $T$ on $V$, and the matrix $E$ used to constrain the frames.

Define \emph{linear testers} on moments by selecting ordered matrix entries:
\begin{align*}
\eta_{d,t}(w(s,z))&=\sum_{\rho=1}^{r_0}z_{O_{d,t},2\rho-1}z_{O_{d,t},2\rho},\,d,t\in I\\
\eta_S(w(s,z))&=\sum_{\rho=1}^{r_0}z_{S,2\rho-1}z_{S,2\rho}.
\end{align*}
Put $\eta_O=\sum_{d,t}\eta_{d,t}$ and $\eta=\eta_O+\eta_S$. These are elements of $\mathcal M^*$ and take values in $\mathbb F_2$. 

For a cut profile $x$ represented by $(w_\ell)$, set
\[
a_t(x)=\sum_{\ell,d}\eta_{d,t}(w_\ell),\qquad
a(x)=\sum_t a_t(x),\qquad
b_t(x)=\sum_{\ell\ne t}\eta_S(w_\ell).
\]
These functionals do not depend on the representation: the $O$-testers vanish on $\mathcal M_*$, and a constant change contributes $(g-1)\eta_S(w_*)=0$ to $b_t$.

Define
\[
T_0(x,y)=a(x)a(y)+\sum_{t\in I}
\bigl(a_t(x)b_t(y)+b_t(x)a_t(y)\bigr).
\]
For each $1\le j\le J$ and ordered pair $e,f\in\mathcal E$, fix endomorphisms $L_j^{ef},R_j^{ef}$ of $B$, chosen to satisfy the uniform mixing conditions of Lemma~5.5 of the manuscript. With the entrywise pairing $\langle A,B\rangle=\operatorname{Tr}(A^TB)$, put
\begin{align*}
B_1(x,y)&=\sum_{j=1}^J\sum_{e,f\in\mathcal E}
\left\langle x_e,L_j^{ef}y_f(R_j^{ef})^T\right\rangle,\\
T_1(x,y)&=B_1(x,y)+B_1(y,x),\qquad T=T_0+T_1.
\end{align*}
Then $T$ is symmetric and $T(x,x)=a(x)$. This specifies $r(\lambda)=a+T(\lambda,\cdot)$ in the edge definition.

The matrix $E\in\operatorname{Mat}_{d_B}(\mathbb F_2)$ prescribes $P_{o,e}^TQ_{o,e}=E$ in Section~\ref{sec:frames}. Let $E_\eta$ be the sum of the ordered tester matrices, so $\langle E_\eta,w\rangle=\eta(w)$ for $w\in\mathcal M$. Choose $M_1,\ldots,M_b\in\operatorname{Mat}_n(\mathbb F_2)$ along with the mixers as in Lemma~5.5, and define $E_{\#}$ by
\[
v(s,z)^TE_\#v(s',z')=
\sum_{\rho=1}^b(s_\rho+s'_\rho)z_\#^TM_\rho z'_\#.
\]
This is bilinear on $B$ since $j_*\ge1$. Set $E=E_\eta+E_\#$. Its additional term vanishes on coincident point inputs, so
\[
\langle E,w\rangle=\eta(w)\qquad(w\in\mathcal M).
\]
The $Z$-block does not occur in $E$. The mixers and the matrices $M_\rho$ are chosen once for each sufficiently large $n$, before graph sampling. Their existence with the required rank properties is supplied by Lemma~5.5; arbitrary choices need not give the connected-matching bound.

\subsection{The vertex law}
With $\mathcal E,d_B,N,h,V,a,T,E$ now specified, apply Section~\ref{sec:frames}. Let $\mathcal O_n$ be the set of admissible frame tuples and let $\mu_n$ be its product uniform law across components. An element $o\in\mathcal O_n$ determines the maps $U_o:V\to W$ and $u_o\in W^*$ already defined there. In the general scheme of Section~\ref{sec:2.1}, take $\Omega=\mathcal O_n$.

The initial random variable is therefore a tuple of frame matrices. The coefficient space, moment spaces, $a,T$, and $E$ are fixed before drawing these tuples. In particular, the matrix $E$ alone does not specify the graph: adjacency also depends on $V,a,T$.

\subsection{The sampled graph and the remaining theorem}
Sample
\[
o_1,\ldots,o_m\ \overset{\mathrm{iid}}{\sim}\ \mu_n,
\qquad m=2^{C_0gN}.
\]
On the positions $1,\ldots,m$, put $ij\in E(G)$ exactly when $o_i,o_j$ satisfy the edge criterion~\eqref{eq:shared}--\eqref{eq:parity}; put $ij\in E(H)$ otherwise. This completes the definition of the random graph. Its edges are determined by the sampled frames, with no further independent edge sampling.

\begin{remark}
The initial random variable is the tuple of frame matrices $o$, not a Boolean point $(s,z)$. The monomial list and moment spaces are fixed globally. Boolean points vary when constructing witnesses; sampling point parameters with the frames fixed is called a query in the manuscript.
\end{remark}

Triangle-freeness of $G$, and hence $\alpha(H)\le2$, holds for every sample. The harder assertion is that $\cm(H)<m/100$ with probability $1-\exp(-\Omega(m))$ (Proposition~3.4). To obtain it, the manuscript calls an ordered pair of raw vertices that do not satisfy the edge criterion a \emph{unit}; such a pair represents an edge of $H$ when placed at distinct sampled positions. Two units conflict if all four cross pairs are edges of $G$, so no edge of $H$ joins their endpoint sets. Their corresponding edges therefore cannot coexist in a connected matching.

Theorem~3.1 asserts that every law $\sigma$ on units satisfying the pointwise density bounds
\[
\sigma_1\le M\mu_n,\qquad \sigma_2\le M\mu_n,\qquad
\sigma\le2^{DN}(\mu_n\otimes\mu_n)
\]
gives a conflict between two independent $\sigma$-units with probability at least $2^{-100gN}$. The endpoints within one unit may be dependent. An entropy-based container argument transfers this statement to the sampled graph. Thus the construction above specifies the graph; the subsequent distribution theorem and container argument establish the small connected matching number needed for the counterexample.
\section{From Theorem 3.1 to the counterexample}
\label{sec:containers}
Assume Theorem~3.1 of the manuscript. We explain how it implies that the sampled graph $H=\bar G$ has $\cm(H)<m/100$ with probability $1-\exp(-\Omega(m))$. Throughout this section write $\Omega=\mathcal O_n$, $\mu=\mu_n$, and $\mu^2=\mu\otimes\mu$.

\subsection{The distribution theorem and the conflict graph}
Let $\mathcal U\subseteq\Omega^2$ be the set of units: ordered pairs that do not satisfy the edge criterion for $G$. Define a conflict graph $\Gamma$ on $\mathcal U$, joining $(a_1,a_2)$ to $(b_1,b_2)$ when all four cross pairs $a_ib_j$ satisfy that criterion. This graph has no loops, since a unit cannot conflict with itself. A connected matching of $H$, after orienting its edges, gives an independent set of unit types in $\Gamma$. Different matching edges may have the same type; we keep track of their occurrences in the sample separately.

For a probability law $\sigma$ on units, its marginals are
\[
\sigma_1(a)=\sum_b\sigma(a,b),\qquad
\sigma_2(b)=\sum_a\sigma(a,b).
\]
Call $\sigma$ feasible if the following inequalities hold pointwise:
\begin{equation}
\sigma_1\le M\mu,\qquad \sigma_2\le M\mu,\qquad
\sigma\le2^{DN}\mu^2.
\label{eq:feasible-law}
\end{equation}
These are upper bounds on concentration, rather than distances between probability laws. Theorem~3.1 says that two independent draws from any feasible law conflict with probability at least
\[
\varepsilon=2^{-100gN}.
\]
The law need not be uniform, and its two endpoints need not be independent.

\subsection{Containers obtained by deleting conflicts}
Fix an independent set $I\subseteq\mathcal U$. Start with the residual family $R=\mathcal U$. Whenever $R$ supports a feasible law, choose one, denoted $\rho_R$, minimizing relative entropy with respect to $\mu^2$:
\[
\mathcal D(\rho_R\Vert\mu^2)
=\sum_{a,b}\rho_R(a,b)\log\frac{\rho_R(a,b)}{\mu(a)\mu(b)}.
\]
Here logarithms are natural and zero summands are omitted. By Theorem~3.1,
\[
\sum_{v\in R}\rho_R(v)\rho_R(N_R(v))\ge\varepsilon,
\]
where $N_R(v)$ is the conflict neighborhood of $v$ in $R$. Choose $v$ maximizing $\rho_R(N_R(v))$, so this mass is at least $\varepsilon$. Make all choices deterministic, using a fixed ordering to break ties.

If $v\in I$, record $v$ and delete $N_R(v)$; if $v\notin I$, delete $v$ alone. Both steps preserve $I\subseteq R$. Stop when no feasible law remains. The recorded units form a fingerprint, and the terminal family $R$ is its container.

The entropy bound limits the number of recordings. Every feasible law has
\[
0\le\mathcal D(\rho\Vert\mu^2)\le DN\log2.
\]
The information-projection inequality of Lemma~3.2 says that, after a recording which leaves a feasible residual family, the minimum entropy increases by at least $-\log(1-\varepsilon)$. Indeed, the old minimizer assigns mass at most $1-\varepsilon$ to the new residual family, and the new law is supported there. Other deletions cannot decrease the minimum entropy. Allowing one final recording that destroys feasibility, the number of recordings is at most
\[
L_N=1+\left\lceil\frac{DN\log2}{-\log(1-\varepsilon)}\right\rceil
=O(N2^{100gN}).
\]
The fingerprint determines the terminal family by replaying the deterministic procedure: a selected unit is recorded exactly when it belongs to the fingerprint. A previously recorded unit has no remaining conflict neighbors and cannot be selected again while feasibility holds.

There are at most $2^{2|\mathcal E|N(d_B+h)}$ raw vertices, so $\log|\mathcal U|=O(N^2)$. Hence the number of containers is at most $(|\mathcal U|+1)^{L_N}$, whose logarithm is
\[
O(N^3 2^{100gN})=o(m),\qquad m=2^{1000gN}.
\]
Thus every independent set lies in one of $\exp(o(m))$ terminal containers, all determined before the graph sample is drawn.

\subsection{A terminal container has rare endpoints or rare pairs}
If a family $R$ supports no feasible law, Lemma~3.3 gives sets $S\subseteq\Omega$ and $E_0\subseteq\Omega^2$ such that
\[
\mu(S)<1/M,\qquad \mu^2(E_0)<2^{-DN},
\]
and each $(a,b)\in R$ has $a\in S$, $b\in S$, or $(a,b)\in E_0$.

To see why, put a left and a right copy of $\Omega$ in a flow network. Give the edges from the source to left $a$ capacity $M\mu(a)$, those from right $b$ to the sink capacity $M\mu(b)$, and the edge from left $a$ to right $b$ capacity $2^{DN}\mu(a)\mu(b)$ whenever $(a,b)\in R$. A flow of value one is precisely a feasible law. If none exists, a cut of capacity less than one supplies sets $L_0,R_0\subseteq\Omega$ with
\[
M\mu(L_0)+M\mu(R_0)
+2^{DN}\mu^2\bigl(R\cap((\Omega\setminus L_0)\times(\Omega\setminus R_0))\bigr)<1.
\]
Take $S=L_0\cup R_0$ and let $E_0$ be the pair set in the last term. This proves the asserted cover and measure bounds.

\subsection{The random sample excludes large connected matchings}
Fix one terminal container and its sets $S,E_0$, and put $k=\lceil m/200\rceil$. For independent sampled positions $o_1,\ldots,o_m$, the probability that at least $k$ lie in $S$ is at most
\[
\binom{m}{k}\mu(S)^k
\le\left(\frac{\mathrm e m}{kM}\right)^k
\le(200\mathrm e/M)^k
=\exp(-\Omega(m)).
\]
Here $\mathrm e$ is Euler's number. A fixed collection of $k$ disjoint ordered position pairs consists of independent $\mu^2$ draws. The probability they all belong to $E_0$ is at most $2^{-DNk}$. There are at most $m^{2k}$ such collections, so
\[
\Pr(\text{$k$ disjoint sampled pairs lie in }E_0)
\le m^{2k}2^{-DNk}
=2^{(2C_0g-D)Nk}
=2^{-2C_0gNk}.
\]
Taking a union bound over the $\exp(o(m))$ containers, with probability $1-\exp(-\Omega(m))$ every container has fewer than $m/200$ sampled positions in $S$ and fewer than $m/200$ disjoint sampled pairs in $E_0$.

Suppose on this event that $H$ has a connected matching of size at least $m/100$. Orient its edges. Their unit types form an independent set in $\Gamma$, so they belong to a terminal container. Because the matching uses disjoint positions, fewer than $m/200$ of its edges can meet that container's sampled positions in $S$. All remaining edges belong to $E_0$, and there are fewer than $m/200$ of those as well. This contradicts the assumed matching size. Therefore
\[
\cm(H)<m/100.
\]
No assumption that the uniform law on the matching itself is feasible is needed: the container procedure applies to its independent set of types regardless.

Finally, $\alpha(H)\le2$ holds for every sample. Applying Appendix~\ref{app:minor-bound}, for $m\ge5$ we obtain
\[
h(H)\le\frac{m+4\cm(H)+2}{3}
<\frac{26m}{75}+\frac23
<\frac m2\le\chi(H).
\]
Thus Theorem~3.1 implies that the construction produces a counterexample with probability $1-\exp(-\Omega(m))$ for all sufficiently large $n$.

\appendix
\section{The clique-minor bound (Proposition 3.5)}
\label{app:minor-bound}

\begin{proposition}
For every finite nonempty graph $H$ with $m$ vertices,
\[
h(H)\le \frac{m+4\cm(H)+2}{3}.
\]
\end{proposition}

\begin{proof}
Write $c=\cm(H)$ and $b=h(H)$. A $K_b$ minor model consists of $b$ disjoint nonempty connected branch sets, with an edge joining every pair. Let $s$ be the number of singleton branch sets and $e$ the number of two-vertex branch sets.

Each two-vertex branch set supplies an edge. The singleton vertices form a clique, so pairing them supplies $\lfloor s/2\rfloor$ more edges. Together these form a connected matching, since every pair of its edges is joined by an edge of $H$. Thus
\[
c\ge e+\left\lfloor\frac{s}{2}\right\rfloor,
\qquad s\le2(c-e)+1.
\]
Every remaining branch set has at least three vertices, so
\[
m\ge s+2e+3(b-s-e)=3b-2s-e.
\]
Consequently,
\[
3b\le m+2s+e\le m+4c-3e+2\le m+4c+2,
\]
which proves the claim.
\end{proof}
\end{document}
